\documentclass{article}
\usepackage{amsmath, amsfonts, amssymb}
\usepackage{graphicx}
\usepackage{titlesec}
\usepackage{lipsum}
\usepackage[utf8]{inputenc}
\usepackage{xcolor}
\usepackage{listings}
\usepackage{lipsum}
\usepackage{scrextend}
\usepackage{tikz}
\usepackage{hyperref}
\usepackage{color}
\usepackage{fancyvrb}
\usepackage[left=2.5cm,top=3cm,right=2.5cm,bottom=3cm,bindingoffset=0.5cm]{geometry}

\title{Adversarial Constraint Calculus}
\author{Luc Alexander L{\"u}thi}
\date{April 10th, 2026}

\begin{document}
	\maketitle
	\tableofcontents
	\section{Introduction}
	\section{Formal Definition}
	A proof consists of an object $p$ such that
	$$p = (\Gamma, \Sigma)$$
	where $\Gamma$ and $\Sigma$ are sets of persistent proof context and mutable state respectively such that
	$$\Gamma \subseteq \Gamma'$$
	$$\Sigma \rightarrow \Sigma'$$
	We can assert that a proof $p$ shows $\phi$ as $p \vdash \phi$ or more formally $\Gamma \vdash \phi$ for some $p=(\Gamma,\Sigma)$. Similarly we can assert that a proof $p$ models $\phi$ as $p \models \phi$ iff more formally $\Sigma \vdash \phi$ for some $p=(\Gamma,\Sigma)$. A property $\phi$ is sound under a proof $p$ iff $p\vdash \phi \wedge p\models \phi$.\\\\
	Instructions are total or partial functions $f : (\Gamma,\Sigma) \rightarrow (\Gamma',\Sigma')$ subject to $\Gamma'=\Gamma\cup\Delta_f(\Gamma,\Sigma)$. These act as parametric theorems which may construct a proof object from another proof object, or may return another parametric theorem.\\\\
	Failure of any instruction or assertion in a proof fails the entire proof. We represent failure of a proof or assertion as $\bot$ such that
	\[
	f(p) =
	\begin{cases}
		(\Gamma', \Sigma') & \text{if preconditions are satisfied} \\
		\bot & \text{otherwise}
	\end{cases}
	\]
	A construction rule constitutes addition of information to the context of a proof.
	$$\text{construct}_\phi(\Gamma,\Sigma) = (\Gamma\cup{\phi},\Sigma)$$
	We may similarly apply a state transition to a proof state.
	$$\text{update}_u(\Gamma,\Sigma) = (\Gamma,u(\Sigma))$$
	Where $u:\Sigma\rightarrow\Sigma$. We notate conditional construction construction of $b$ if $a$ as an implication $a \rightarrow b$ and conditional update on of $b$ if $a$ as an implication $a \Rightarrow b$.
	Naturally $(f \circ g)(p) = f(g(p))$.
	\section{Operational Semantics}
	An instruction f constitutes theorem application on $p$, and undergoes a transition relation
	$$(\Gamma,\Sigma)\xrightarrow{f}(\Gamma',\Sigma') \text{ iff } f(\Gamma,\Sigma) = (\Gamma',\Sigma')$$
	Execution is traced for $p_0\xrightarrow{f_1}p_1\xrightarrow{f_2}\text{...}\xrightarrow{f_n}p_{n}$ such that
	$$\Gamma_0\subseteq\Gamma_1\subseteq\text{...}\subseteq\Gamma_{n}$$
	An invariant is a property over a state $I(\Sigma)$, and is preserved if $\forall f, I(\Sigma) \xrightarrow{f} I(\Sigma')$. An invariant is certified if $\Gamma \vdash I$, and is sound only when it is both shown and modeled.
	\section{Adversarial Modeling}


\end{document}
